% Generated by tex/build_swebok_structure.py from tex/chapters/ch01/05_要求分析.tex
\subsubsection{要求における衝突への対応}

利害関係者が多いほど、要求の衝突は起こりやすい。たとえば、利用者は高機能を求め、運用担当は単純で安定した運用を求め、経営側は費用や納期を抑えたいと考えることがある。

衝突への対応には、主に二つの方向がある。第一に、関係者間で交渉し、合意を形成することである。技術者が独断で決めると、責任や契約上の問題が残ることがある。第二に、\term{製品群開発}の考え方を使い、すべての利用者に共通する不変要求と、利用者や市場によって変わる可変要求を分けることである。

\MiniBox{例}{天気アプリで、ある利用者は摂氏表示を求め、別の利用者は華氏表示を求める。この場合、温度を表示することは共通要求であり、単位は可変要求である。}
